\documentclass[11pt,a4paper]{article}
\usepackage{fontspec}
\setmainfont{DejaVu Sans}[Scale=0.90]
\usepackage{amsmath,amssymb}
\usepackage[margin=2.2cm]{geometry}
\usepackage{booktabs,array,tabularx}
\usepackage{enumitem}
\usepackage{titlesec}
\usepackage{parskip}
\titlelabel{\thetitle.\quad}
\titleformat{\section}{\normalfont\large\bfseries}{\thesection.}{0.6em}{}
\titleformat{\subsection}{\normalfont\normalsize\bfseries}{\thesubsection}{0.6em}{}
\titlespacing*{\section}{0pt}{14pt}{5pt}
\titlespacing*{\subsection}{0pt}{10pt}{4pt}
\setlist[itemize]{leftmargin=1.4em,itemsep=1pt,topsep=3pt}

% Preserve fonts/margins; permit a final line-breaking pass for long prose.
\setlength{\emergencystretch}{2em}
\begin{document}

\begin{center}
{\Large\bfseries Admissibilité des lois de conduction non-Fourier}\\[5pt]
{\normalsize Note de partie II --- point de contrôle}\\[2pt]
{\small Version corrigée --- 15 septembre 2026}
\end{center}

\vspace{4pt}

\section{Objet}

Deux critères sont couramment invoqués pour juger une loi de conduction : la compatibilité avec le
second principe, et la finitude de la vitesse de propagation. Le second est le grief habituel contre
la loi de Fourier, et le motif usuel d'introduire un temps de relaxation.

Ces deux critères sont fréquemment présentés comme allant de pair, la hiérarchie Fourier,
Cattaneo, Guyer--Krumhansl étant décrite comme un raffinement monotone. La présente note établit
qu'ils sont indépendants, et que cette hiérarchie n'est pas monotone.

\section{Résultat}

\begin{center}
\begin{tabularx}{\textwidth}{@{}lXXl@{}}
\toprule
\textbf{Loi} & \textbf{Entropie math.} & \textbf{Second principe} & \textbf{Propagation} \\
\midrule
Fourier & oui & oui, si $\lambda > 0$ & infinie \\
\addlinespace
Cattaneo & oui, si $\tau_R > 0$ & oui, si $\lambda > 0$ & finie, $\sqrt{a/\tau_R}$ \\
\addlinespace
Guyer--Krumhansl & oui, si $\tau_R > 0$ & oui, si $\lambda > 0$ et $\ell^2 \ge 0$ & infinie si $\ell>0$ \\
\bottomrule
\end{tabularx}
\end{center}

\medskip
Le tableau emploie l'opposé de l'entropie physique, qui est concave. Les trois lois sont admissibles près de l'équilibre. Cattaneo propage à vitesse finie ; GK avec $\ell=0$ se réduit à ce cas. Le
modèle le plus riche des trois perd la propriété qui motive habituellement l'abandon de Fourier.

\section{Vitesse de propagation}

\subsection{Relation de dispersion}

La combinaison de la loi constitutive et du bilan d'énergie donne, pour une onde plane en milieu
homogène,
\[ \sigma^2 = \frac{p\,(1+\tau_R p)}{a\,(1+\tau_\ell p)} , \qquad p = i\omega , \qquad \sigma = -ik , \]
où le champ est proportionnel à $\exp(i\omega t+ikz)$ et où $\tau_\ell = 3\ell^2/a$ est le temps de diffusion sur la longueur non locale. Cattaneo correspond
à $\tau_\ell = 0$, Fourier à $\tau_R = \tau_\ell = 0$.

Le comportement du nombre d'onde à haute fréquence décide de la vitesse. Les exposants mesurés sur
quatre décades :

\begin{center}
\begin{tabular}{@{}lcl@{}}
\toprule
\textbf{Loi} & $k \propto \omega^{\,n}$ & \textbf{Vitesse de phase} \\
\midrule
Fourier & $n = 0{,}5000$ & non bornée \\
Cattaneo & $n = 1{,}0000$ & sature à $\sqrt{a/\tau_R}$ \\
Guyer--Krumhansl ($\ell>0$) & $n = 0{,}5004$ & non bornée \\
\bottomrule
\end{tabular}
\end{center}

\subsection{Interprétation}

Un exposant $1$ signifie une vitesse de phase constante : une onde véritable. Un exposant $1/2$ est
la signature de la diffusion, et la vitesse de phase croît alors comme la racine de la fréquence,
sans borne.

\textbf{Le terme non local est diffusif sur le flux lui-même.} L'équation de Guyer--Krumhansl
comporte un terme en $\partial^2 q / \partial z^2$, du second ordre en espace et du premier en temps
pour le flux. Au-dessus de la fréquence où ce terme domine celui en $\tau_R\,\partial_t q$, la partie
principale du système redevient parabolique, et la vitesse finie que Cattaneo avait restaurée est
perdue.

\subsection{Système caractéristique de Cattaneo}

En variables $(u, q)$, la partie principale s'écrit $\partial_t U + A\,\partial_z U = 0$ avec
\[ A = \begin{pmatrix} 0 & 1 \\[2pt] a/\tau_R & 0 \end{pmatrix} ,
\qquad \text{valeurs propres } \pm\sqrt{a/\tau_R} . \]

Deux valeurs propres réelles distinctes : le système est strictement hyperbolique. Une entropie mathématique strictement convexe peut symétriser un système de conservation sous les hypothèses du théorème de Godunov--Mock ; l'hyperbolicité seule ne prouve pas l'existence d'une telle entropie. La section suivante vérifie directement l'entropie de ce modèle.

En modèle de Debye, $a = v^2\tau_R/3$, de sorte que la vitesse limite vaut $v/\sqrt{3}$ : elle est
fixée par la vitesse du son et ne contient aucune information supplémentaire.

\section{Entropie}

\subsection{Cadre}

La thermodynamique irréversible étendue prend le flux comme variable d'état à part entière :
\[ s(u, q) = s_{\mathrm{éq}}(u) - \frac{\tau_R}{2\lambda T_0^{\,2}}\;q^2 . \]

L'entropie physique est concave ; son opposé, utilisé comme entropie mathématique,
est convexe. La partie hors équilibre est une forme quadratique définie négative dès que $\tau_R > 0$ et
$\lambda > 0$ : l'entropie est concave, donc l'état d'équilibre est stable.

\subsection{Productions}

Avec le flux classique pour Cattaneo et le flux étendu ci-dessous pour GK,
la dérivation depuis les équations de bilan
donne, à l'ordre quadratique des perturbations autour d'une température de référence $T_0$,
\[ \boxed{\;\sigma_s^{\text{Cat}} = \frac{q^2}{\lambda T_0^{\,2}}\;}
\qquad\qquad
\boxed{\;\sigma_s^{\text{GK}} = \frac{q^2 + 3\ell^2\,(\partial_z q)^2}{\lambda T_0^{\,2}}\;} \]

Les deux sont des sommes de carrés à coefficients positifs. Le second principe est donc satisfait
dans les deux cas, sous les seules conditions $\lambda > 0$ et $\ell^2 \ge 0$.

\subsection{Le flux d'entropie doit être étendu}

Le second résultat n'est pas obtenu avec le flux classique. Il exige
\[ \mathbf{J}_s = \frac{\mathbf{q}}{T} + \frac{3\ell^2}{\lambda T_0^{\,2}}\;q\,\partial_z q , \]
le coefficient étant déterminé, et non postulé, par l'annulation du terme croisé en
$q\,\partial_z^2 q$. Sans ce terme supplémentaire, la production conserve une contribution de signe
indéfini et n'est pas une somme de carrés.

\subsection{Une erreur de dérivation à signaler}

Le flux d'entropie doit être pris en $\mathbf{q}/T$ et non en $\mathbf{q}/T_0$, même lorsque la
linéarisation est appliquée partout ailleurs. La différence entre les deux est exactement le terme
qui compense la contribution du gradient de température ; avec $\mathbf{q}/T_0$ il subsiste un
résidu en $q\,\partial_z T$ qui ne s'annule pas et qui conduit à conclure à tort que la production
n'est pas définie positive.

\section{Portée et limites}

\subsection{Conséquence pour le choix d'un modèle}

Le terme non local de Guyer--Krumhansl produit trois effets distincts, établis séparément :

\begin{itemize}
\item il \textbf{ajoute} un paramètre, $\tau_\ell$, dont l'identifiabilité dépend
de la bande, du bruit et des corrélations ; le libérer peut dégrader les autres estimations ;
\item il \textbf{ne lève pas} la dégénérescence d'échelle du problème inverse, n'entrant dans
l'observable que par un temps ;
\item pour $\ell>0$, il \textbf{détruit} la finitude de la vitesse de propagation ; $\ell=0$ redonne Cattaneo.
\end{itemize}

Une liberté supplémentaire, une invariance conservée et un changement de propagation. Le choix entre Cattaneo
et Guyer--Krumhansl ne se réduit donc pas à retenir le modèle le plus complet.

\subsection{Ce qui n'est pas établi}

Le modèle à deux temps de relaxation, dit à retard de phase double, n'est pas traité ici. Son
caractère mal posé au-delà du premier ordre de développement est établi dans la littérature ; la
conséquence pour les valeurs de propriétés extraites en métrologie de couches minces reste à
examiner.

Les dérivations sont menées en une dimension et à propriétés constantes. Le cas gradué sous loi
non-Fourier n'est pas abordé.

L'admissibilité est établie au sens de la thermodynamique irréversible étendue, qui postule la forme
de l'entropie hors équilibre. D'autres cadres --- thermodynamique rationnelle étendue, formalisme
lagrangien --- conduisent aux mêmes équations par d'autres voies ; leur comparaison n'est pas menée.

Enfin, la vitesse infinie de Guyer--Krumhansl est une propriété du modèle continu. Elle n'a de portée
physique que dans le domaine de validité de celui-ci, qui suppose des variations lentes devant le
libre parcours moyen. Le paradoxe est donc atténué par le fait que le modèle ne prétend pas décrire
les fréquences où il se manifeste.

\section{Références}

{\small
\begin{enumerate}[leftmargin=1.6em,itemsep=0pt]
\item C. Cattaneo, \emph{Sulla conduzione del calore}, Atti Sem. Mat. Fis. Univ. Modena \textbf{3}, 83 (1948).
\item R. A. Guyer, J. A. Krumhansl, Phys. Rev. \textbf{148}, 766 et 778 (1966).
\item D. Jou, J. Casas-Vázquez, G. Lebon, \emph{Extended Irreversible Thermodynamics}, Springer.
\item S. K. Godunov, \emph{An interesting class of quasilinear systems}, Dokl. Akad. Nauk SSSR
\textbf{139}, 521 (1961).
\item M. S. Mock, \emph{Systems of conservation laws of mixed type}, J. Differential Equations
\textbf{37}, 70 (1980).
\end{enumerate}}

\vspace{4pt}
\noindent\small
Vérifications : \texttt{tests/\allowbreak test\_\allowbreak admissibility.py}, dont la dérivation symbolique des deux
productions d'entropie. Implantation : \texttt{src/\allowbreak admissibility.py}. Conventions :
\texttt{notes/\allowbreak conventions.md}, section~10. Références à vérifier à la source avant citation.

\end{document}
